%% =====================================================================
%%  membretes/despacho/membrete.tex  --  Variante Camara de Diputados
%% ---------------------------------------------------------------------
%%  Version moderna del membrete usando el logotipo oficial de la
%%  CAMARA DE DIPUTADOS (bicameralidad 2026) junto al escudo nacional.
%%  Mismo contrato que membretes/congreso: nombre/cargo en Selawik,
%%  lema centrado respecto a la caja y en tono sutil.
%%
%%  Seleccionar con:  \usarmembrete{despacho}
%% =====================================================================

\fancyhead[C]{%
  \begin{minipage}[b]{\textwidth}%
    % --- Fila unica: logos a la izquierda | nombre + lema a la derecha
    \begin{minipage}[c]{0.44\textwidth}%
      \raggedright
      \includegraphics[height=1.50cm]{logo-camara-diputados.png}%
      \hspace{0.30cm}%
      \includegraphics[height=1.42cm]{escudo-peru.png}%
    \end{minipage}%
    \hfill
    \begin{minipage}[c]{0.54\textwidth}%
      \raggedleft
      \setlength{\fboxsep}{0pt}%
      \colorbox{colormembrete}{%
        \begin{minipage}[c][1.30cm][c]{\despachoAnchoCaja}%
          \centering\color{white}%
          {\fuentemembretesemi\fontsize{10.5}{13}\selectfont
            \MakeUppercase{\meta{autor}}}%
          % Cargo vacio -> caja de UNA linea (p. ej. «BANCADA REGIONAL»,
          % \setmeta{cargo}{} en su metadata; pedido Edison 2026-08-23).
          \ifmetadef{cargo}{\\[2.5pt]
          {\fuentemembrete\fontsize{9}{11}\selectfont
            \MakeUppercase{\meta{cargo}}}}{}%
        \end{minipage}%
      }\\[4.5pt]
      % --- Lema del anio: centrado respecto a la caja, tono sutil ----
      \makebox[\despachoAnchoCaja][c]{%
        {\itshape\fontsize{7.5}{9}\selectfont\color{neutral700}%
          ``\meta{lema}''}%
      }%
    \end{minipage}%
    \par\vspace{5pt}
    % --- Filete ------------------------------------------------------
    \ifdespachofileteencabezado{\color{colorfilete}\rule{\textwidth}{0.6pt}}\fi%
  \end{minipage}%
}
